\documentclass[10pt]{article}
\usepackage[utf8]{inputenc}
\usepackage[T1]{fontenc}
\usepackage[brazil]{babel}
\usepackage[a4paper,margin=2cm]{geometry}
\usepackage{booktabs}
\usepackage{array}
\usepackage{amsmath}
\usepackage{xcolor}
\usepackage{hyperref}
\hypersetup{colorlinks=true,urlcolor=blue!60!black,linkcolor=black}

\newcommand{\sd}{\ensuremath{S^{\delta}}}
\setlength{\parindent}{0pt}
\setlength{\parskip}{3pt}

\begin{document}

% ================= CABECALHO COMPACTO =================
\begin{center}
  {\large\bfseries Trabalho 3 (T3): Divisão e conquista com MPI}\\[2pt]
  {\small Ordenação de um vetor (Bubble Sort nas folhas) --- Unidade 04}\\[4pt]
  {\footnotesize
   \textbf{Integrante:} Bernardo Fioreze \quad|\quad
   \textbf{Repositório:} \url{https://github.com/BernardoFioreze/t3-mpi-divisao-conquista}}\\[2pt]
  {\footnotesize
   \textbf{Ambiente:} cluster Atlântica (LAD), \textbf{2 nós exclusivos}; cada nó = 2$\times$ Intel Xeon
   E5520 @ 2,27\,GHz = \textbf{8 núcleos físicos / 16 threads (HT) por nó} (\texttt{lscpu}),
   logo \textbf{16 núcleos físicos / 32 threads} nos 2 nós. $N=200000$, 10 repetições por configuração
   (mediana\,/\,mín\,/\,máx), todas com \texttt{verificacao=OK}.}
\end{center}
\hrule
\vspace{4pt}

% ================= ANALISE =================
\textbf{\large Análise dos resultados}

\textbf{1. Modelagem.}
Todos os processos rodam o mesmo executável (SPMD) e descobrem seu papel pelo \texttt{rank}.
Na \emph{versão original} a árvore é indexada como um heap: a raiz é o \texttt{rank}~0, um processo
$r$ tem filhos $2r{+}1$ e $2r{+}2$ e pai $(r{-}1)/2$ (\texttt{original.c}:9--10). O processo não conhece
o tamanho do pedaço de antemão: ele o descobre com \texttt{MPI\_Probe}/\texttt{MPI\_Get\_count} antes de
alocar (\texttt{original.c}:99--102). Se o pedaço é $\le\delta$ ele conquista com Bubble Sort
(linhas 123--124); senão envia \textbf{as duas metades} aos filhos (linhas 129--130) \textbf{antes} de
qualquer \texttt{MPI\_Recv} (linhas 132--133), para os dois filhos trabalharem ao mesmo tempo, e
intercala na volta (linha 134). Na \emph{versão balanceada} o processo não entrega as duas metades: manda
a metade direita ao ajudante $r+2^{\text{nível}}$ (\texttt{MPI\_Isend}, \texttt{balanceado.c}:90) e
\textbf{fica com a esquerda}, descendo um nível (linha 95); recebe a direita ordenada de volta e intercala
(linhas 98--100). Um não-raiz descobre onde nasceu pelo bit menos significativo do \texttt{rank}:
nível inicial $=\texttt{ctz}(r)-1$ (linha 170) e pai $= r\,\&\,(r{-}1)$ (linha 139). O \texttt{MPI\_Isend}
seguido de \texttt{MPI\_Request\_free} (linhas 90--92) evita bloquear e descarta a \emph{request} porque o
\texttt{MPI\_Recv} do mesmo ajudante (linha 98) só retorna depois que a mensagem foi consumida. Com $P$
potência de 2 todas as folhas têm $N/P$; com $P=2^{k}-1$ (caso da original) elas diferem em um elemento.

\textbf{2. Speed-up, eficiência e a separação do ganho de algoritmo.}
A Tabela~\ref{tab:seq} mostra o ganho de algoritmo $T_1/T_1^{\delta}$ indo de $2{,}0\times$ ($\delta{=}10^5$)
a $535\times$ ($\delta{=}391$): quase todo o ``speed-up'' acima de $P$ vem daí, não do paralelismo.
Por exemplo, na original com $P{=}31$ e $\delta{=}12500$ temos $S=125\times$ mas \sd${}=7{,}4\times$; e
$7{,}4 \times 16{,}9$ (o ganho de algoritmo em $\delta{=}12500$) $\approx 125$. A \sd{} é o paralelismo
honesto (contra a \emph{mesma} árvore em um processo) e a eficiência $E=\sd/P$ cai de $0{,}67$ ($P{=}3$) a
$0{,}24$ ($P{=}31$) por causa da comunicação e da árvore (só as folhas calculam --- ponto 3).
Para a versão balanceada com \textbf{31 processos} (não potência de 2), as folhas têm tamanhos diferentes
e um $\delta$ único não reproduz a árvore; por isso ela fica fora da Tabela~\ref{tab:par} (que traz só as
potências de 2, com \sd{} e $E$), mas seu $S$ é válido e medido: $T_P=0{,}205$\,s, $S=197\times$, contra
$125\times$ da original no \emph{mesmo} $P{=}31$ (ponto 5).

\textbf{3. Balanceamento da carga nos níveis (original).}
Cada nó interno divide e então fica \textbf{parado} em \texttt{MPI\_Recv} (\texttt{original.c}:132--133)
até os filhos voltarem; apenas o nível das folhas executa Bubble Sort. Num dado instante os níveis
superiores estão ociosos e, na volta, fazem só a intercalação barata ($O(n)$). É esse desbalanceamento
que a versão balanceada ataca ao fazer o pai também ordenar uma folha.

\textbf{4. Ganho com HT.}
Comparando $P{=}15$ (um processo por núcleo físico; há 16 físicos nos 2 nós) com $P{=}31$ (passa a usar as
threads de Hyper-Threading): o tempo cai de $0{,}709$ para $0{,}322$\,s ($2{,}2\times$), mas o $\delta$
também muda ($25000\!\to\!12500$). Olhando a \sd{}, que isola o paralelismo, o ganho é só
$6{,}93\!\to\!7{,}41$ ($+7\%$). Ou seja, o HT quase não aumenta o paralelismo real: as duas threads de um
núcleo dividem as mesmas unidades de execução e o Bubble Sort é limitado por CPU; o grosso da melhora de
$P{=}15$ para $31$ vem das folhas menores (algoritmo), não de mais processamento simultâneo.

\textbf{5. Ganho com o balanceamento.}
Duas formas, como pede o enunciado. \emph{Sem alterar o código} (mais processos que núcleos --- série de
oversubscription, Tabela~\ref{tab:over}): quem balanceia é o escalonador do SO (ponto 6).
\emph{Alterando o código} (versão balanceada): no mesmo $P{=}31$ o tempo cai de $0{,}322$ para
$0{,}205$\,s ($1{,}57\times$ mais rápido) porque os processos internos deixam de ficar ociosos e também
ordenam uma folha. Em potências de 2 a eficiência da balanceada se mantém alta ($E\approx0{,}49$ até
$P{=}8$) enquanto a da original já está caindo.

\textbf{6. Desempenho de 63 a 1023 processos: espera ativa $\times$ núcleo cedido.}
Com mais processos que as 32 threads de hardware, a \textbf{espera ativa}
(\texttt{OMPI\_MCA\_mpi\_yield\_when\_idle=0}) faz um processo bloqueado em \texttt{MPI\_Recv} ocupar o
núcleo em laço; com \texttt{OMPI\_MCA\_mpi\_yield\_when\_idle=1} ele \textbf{cede o núcleo}. As duas
séries foram medidas com a variável \textbf{definida explicitamente} (0 e 1), como pede a errata do
enunciado, pois deixá-la indefinida pode já ceder o núcleo e fazer as duas medirem a mesma coisa (foi o
que observamos: com a variável indefinida a espera ativa era bem mais rápida, sinal de que o núcleo já
estava sendo cedido). Em oversubscription moderado o \emph{yield} ganha claramente: em $P{=}127$,
$0{,}519$\,s (espera ativa) contra $0{,}184$\,s (cede o núcleo), $2{,}8\times$; em $P{=}63$, $0{,}333$
contra $0{,}214$. A diferença entre as duas séries é exatamente quanto o
escalonador do SO consegue balancear quando os processos cooperam cedendo o núcleo que não estão usando.
A partir de $P{=}255$ os dois pioram (muito além de 32 threads: troca de contexto e pressão de memória
dominam) e a eficiência $E=\sd/32$ fica ínfima, mas o \emph{yield} degrada menos. O caso $P{=}1023$ rodou
completo nos dois modos após elevar \texttt{ulimit -n} a 512000 e usar \texttt{opal\_set\_max\_sys\_limits}
(senão estourava o limite de \emph{pipes} do SO ao encaminhar a saída de cada rank).

\textbf{7. Uso de ferramentas de IA.}
Declaramos o uso do Claude (Anthropic) para: entender o enunciado; derivar a implementação a partir do
pseudocódigo (indexação em heap na original; esquema de ``pente''/Radenski na balanceada); escrever os
scripts \texttt{sbatch} e o \texttt{analise.py}; e depurar erros do cluster (falha do PMIx e o limite de
\emph{pipes} em $P{=}1023$). Revisamos o código linha a linha, sabemos explicar cada chamada MPI, e as
decisões de modelagem e de medição são do grupo.

% ================= TABELAS =================
\begin{table}[htbp]\centering\footnotesize
\caption{Referências sequenciais ($N{=}200000$, 1 processo). A coluna \emph{níveis} é o número de níveis
da árvore (nível~1 = vetor inteiro, sem divisão, $\delta=N$ --- é o próprio $T_1$); cada nível a mais
divide as folhas pela metade. Todas rodam em um só processo. A mesma árvore de nível $n$ corresponde a
$P=2^{\,n}-1$ processos na versão original e $P=2^{\,n-1}$ na balanceada (daí os $\delta$ compartilhados
na Tabela~\ref{tab:par}). O ganho de algoritmo $T_1/T_1^{\delta}$ explica por que
$S$ passa de $P$.}
\label{tab:seq}
\begin{tabular}{l r r r r r}
\toprule
nível & $\delta$ & mediana (s) & mín (s) & máx (s) & $T_1/T_1^{\delta}$\\
\midrule
1 ($T_1$) & 200000 & 40,354 & 40,339 & 40,383 &   1,00\\
2         & 100000 & 20,051 & 20,033 & 20,073 &   2,01\\
3         & 50000  &  9,961 &  9,913 &  9,975 &   4,05\\
4         & 25000  &  4,912 &  4,904 &  4,921 &   8,22\\
5         & 12500  &  2,384 &  2,379 &  2,387 &  16,93\\
6         & 6250   &  1,125 &  1,113 &  1,134 &  35,87\\
7         & 3125   &  0,552 &  0,547 &  0,562 &  73,14\\
8         & 1563   &  0,290 &  0,287 &  0,292 & 139,35\\
9         & 782    &  0,152 &  0,149 &  0,155 & 265,23\\
10        & 391    &  0,075 &  0,075 &  0,081 & 535,04\\
\bottomrule
\end{tabular}
\end{table}

\begin{table}[htbp]\centering\footnotesize
\caption{Original ($P=2^{\,n}{-}1$) e balanceada ($P=2^{\,n-1}$) lado a lado, por nível $n$ da árvore.
No mesmo nível as duas ordenam folhas do \textbf{mesmo tamanho} ($\delta=N/2^{\,n-1}$, coluna única);
o que muda é quantos processos cada uma usa para chegar lá. $S=T_1/T_P$; \sd${}=T_1^{\delta}/T_P$;
$E=\sd/P$; tempos em segundos. O nível~1 (original com $P{=}1$, vetor inteiro) coincide com $T_1$ da
Tabela~\ref{tab:seq} ($T_P{=}40{,}39$, $S{=}1{,}00$); a balanceada começa no nível~2 e a original não
alcança o nível~6.}
\label{tab:par}
\begin{tabular}{c c @{\hskip 1em} r r r r r @{\hskip 1.5em} r r r r r}
\toprule
 & & \multicolumn{5}{c}{\textbf{original}} & \multicolumn{5}{c}{\textbf{balanceada}}\\
\cmidrule(lr){3-7}\cmidrule(lr){8-12}
nível & $\delta$ & $P$ & $T_P$ & $S$ & \sd & $E$ & $P$ & $T_P$ & $S$ & \sd & $E$\\
\midrule
2 & 100000 &  3 & 10,053 &   4,01 & 1,99 & 0,665 &  2 & 20,081 &   2,01 & 1,00 & 0,499\\
3 &  50000 &  7 &  2,577 &  15,66 & 3,87 & 0,552 &  4 &  4,992 &   8,08 & 1,99 & 0,499\\
4 &  25000 & 15 &  0,709 &  56,90 & 6,93 & 0,462 &  8 &  1,262 &  31,98 & 3,90 & 0,487\\
5 &  12500 & 31 &  0,322 & 125,38 & 7,41 & 0,239 & 16 &  0,361 & 111,84 & 6,61 & 0,413\\
6 &   6250 & \multicolumn{5}{c}{--- (original vai só até $P{=}31$) ---} & 32 &  0,170 & 236,72 & 6,56 & 0,205\\
\bottomrule
\end{tabular}
\end{table}

\begin{table}[htbp]\centering\footnotesize
\caption{Oversubscription da versão original (63--1023 processos nos mesmos 2 nós = 32 threads).
Duas execuções, com a variável definida explicitamente: espera ativa
(\texttt{OMPI\_MCA\_mpi\_yield\_when\_idle=0}) e cede o núcleo
(\texttt{=1}). $E=\sd/32$ (eficiência relativa às 32 threads de hardware).}
\label{tab:over}
\begin{tabular}{l r r r r r r @{\hskip 2em} r r r r r}
\toprule
 & \multicolumn{5}{c}{espera ativa (\texttt{yield\_when\_idle=0})} & & \multicolumn{5}{c}{cede o núcleo (\texttt{=1})}\\
\cmidrule(lr){2-6}\cmidrule(lr){8-12}
$P$ & $\delta$ & $T_P$ (s) & $S$ & \sd & $E$ & & $T_P$ (s) & $S$ & \sd & $E$\\
\midrule
  63 & 6250 & 0,333 & 121,2 & 3,38 & 0,106 & & 0,214 & 188,3 & 5,25 & 0,164\\
 127 & 3125 & 0,519 &  77,7 & 1,06 & 0,033 & & 0,184 & 219,8 & 3,01 & 0,094\\
 255 & 1563 & 0,728 &  55,4 & 0,40 & 0,012 & & 0,228 & 177,2 & 1,27 & 0,040\\
 511 &  782 & 1,906 &  21,2 & 0,08 & 0,002 & & 0,427 &  94,5 & 0,36 & 0,011\\
1023 &  391 & 6,534 &   6,2 & 0,01 & 0,000 & & 2,119 &  19,1 & 0,04 & 0,001\\
\bottomrule
\end{tabular}
\end{table}

\vspace{4pt}
\textbf{\large Teste extra 1: $N = 10^6$}

Para ver o comportamento em escala maior, repetimos o experimento com $N{=}10^6$ (3 repetições, tudo num
único job). O foco foi a oversubscription profunda, então o sequencial rodou só os níveis 10--13
(Tabela~\ref{tab:seq1m}; sem $T_1$); por isso original e balanceado aparecem só com o \textbf{tempo bruto}
$T_P$ (sem $S$/\sd). O nível~10 ($\delta{=}1954$, $T_1^{\delta}{=}1{,}74$\,s) é a referência que dá o \sd\
do oversub $P{=}1023$. A
Tabela~\ref{tab:1m} compara as duas versões por nível (mesmo $\delta$ = mesmas folhas). Dois resultados
se confirmam em escala: (i)~no \emph{mesmo} $\delta$ o balanceado é $\approx2\times$ mais lento que o
original (ex.\ nível~2: $474$ vs $237$\,s), pela serialização do topo (o pai calcula antes de o
\texttt{MPI\_Isend} grande progredir, então o ajudante espera); mas no \emph{mesmo} $P$ ele ganha, por
alcançar uma árvore mais profunda. (ii)~A oversubscription tem \textbf{teto} aqui: em 2 nós, só $P{=}1023$
rodou; $P{=}2047$ e $4095$ estouraram o \emph{timeout} de 20\,min na espera ativa e $P{=}8191$ falhou no
\texttt{MPI\_Init} (o runtime não lança 8191 processos). No $P{=}1023$ o \emph{yield} foi
$\mathbf{3{,}1\times}$ mais rápido que a espera ativa ($2{,}25$ vs $7{,}08$\,s).

\begin{table}[htbp]\centering\footnotesize
\caption{$N{=}10^6$, 3 repetições, tempo bruto $T_P$ (sem $T_1$, logo sem $S$/\sd\ para as versões
paralelas). À direita, a única série de oversub que completou ($P{=}1023$), com $E=\sd/32$.}
\label{tab:1m}
\begin{tabular}{c c @{\hskip 1em} r r @{\hskip 1em} r r @{\hskip 2.5em} l r r r}
\toprule
\multicolumn{6}{c}{\textbf{original $\times$ balanceado (tempo $T_P$, s)}} & \multicolumn{4}{c}{\textbf{oversub $P{=}1023$}}\\
\cmidrule(lr){1-6}\cmidrule(lr){7-10}
nível & $\delta$ & orig $P$ & $T_P$ & bal $P$ & $T_P$ & modo & $T_P$ & \sd & $E$\\
\midrule
2 & 500000 &  3 & 237,2 &  2 & 474,2 & espera ativa (=0) & 7,08 & 0,25 & 0,008\\
3 & 250000 &  7 &  60,0 &  4 & 118,7 & cede núcleo (=1)  & 2,25 & 0,77 & 0,024\\
4 & 125000 & 15 &  15,8 &  8 &  30,4 & \multicolumn{4}{l}{\footnotesize $P\ge2047$: \emph{timeout}/falha no lançamento}\\
5 &  62500 & 31 &   6,26 & 16 &   8,20 & & & &\\
6 &  31250 & \multicolumn{1}{c}{---} & --- & 32 &   3,13 & & & &\\
\bottomrule
\end{tabular}
\end{table}

\begin{table}[htbp]\centering\footnotesize
\caption{Referências sequenciais (1 processo) do teste $N{=}10^6$: $T_1^{\delta}$ dos níveis 10--13
(sem $T_1$). O nível~10 ($\delta{=}1954$) é a referência usada para a \sd\ do oversub $P{=}1023$ da
Tabela~\ref{tab:1m}.}
\label{tab:seq1m}
\begin{tabular}{r r r}
\toprule
nível & $\delta$ & $T_1^{\delta}$ (s)\\
\midrule
10 & 1954 & 1,74\\
11 &  977 & 0,91\\
12 &  489 & 0,43\\
13 &  245 & 0,24\\
\bottomrule
\end{tabular}
\end{table}

\vspace{4pt}
\textbf{\large Teste extra 2: oversub com $N = 5\times10^6$}

Isolando a oversubscription num vetor ainda maior (só os tempos dos dois modos, $P{=}63$--$1023$, 3
repetições) aparece um limite do \emph{yield} (Tabela~\ref{tab:5m}). Com \textbf{folhas grandes} ($P$
baixo) cada processo passa quase todo o tempo ordenando (\emph{CPU-bound}) e quase não espera, então os
dois modos \textbf{empatam} ($P{=}63$--$255$). O \emph{yield} só ganha quando a oversubscription aperta
($P\ge511$, 16--32$\times$ as 32 threads): as folhas ficam menores, o peso de espera/intercalação cresce,
e os processos ociosos da espera ativa roubam ciclos. Repare que na espera ativa o tempo \textbf{estaciona}
de $P{=}511$ para $1023$ ($16{,}8\to16{,}6$\,s: mais processos não ajudam), enquanto no \emph{yield}
continua caindo ($10{,}6\to9{,}7$\,s). Comparado ao $N{=}10^6$ (folhas minúsculas, \emph{yield} $3{,}1\times$
em $P{=}1023$), aqui o ganho cai para $1{,}7\times$: o benefício do \emph{yield} \textbf{cresce com a
oversubscription e encolhe quanto mais pesada (\emph{CPU-bound}) é cada folha}.

\begin{table}[htbp]\centering\footnotesize
\caption{Oversub com $N{=}5\times10^6$ (3 repetições, tempo $T_P$ mediano). Razão $=$ ativa$/$yield
(${>}1$ = yield mais rápido).}
\label{tab:5m}
\begin{tabular}{r r r r r}
\toprule
$P$ & $\delta$ & espera ativa $T_P$ (s) & cede núcleo $T_P$ (s) & razão\\
\midrule
  63 & 156250 & 72,7 & 75,7 & 0,96\\
 127 &  78125 & 37,7 & 38,4 & 0,98\\
 255 &  39063 & 19,9 & 20,0 & 0,99\\
 511 &  19532 & 16,8 & 10,6 & 1,58\\
1023 &   9766 & 16,6 &  9,7 & 1,71\\
\bottomrule
\end{tabular}
\end{table}

\vspace{6pt}
\hrule
\vspace{4pt}
{\footnotesize O código completo (\texttt{src/sequencial.c}, \texttt{src/original.c},
\texttt{src/balanceado.c}), os scripts \texttt{sbatch} e as saídas brutas do Slurm estão no repositório
público citado no cabeçalho; as referências de linha acima apontam para esses arquivos.}

\end{document}
